% =============================================================================
% Kapitel 3: Grundlagen
% =============================================================================
% Kompilierung: pdflatex + biber (oder biblatex mit bibtex backend)
% Benötigte Pakete: inputenc, fontenc, babel, amsmath, booktabs,
%                   graphicx, hyperref, biblatex, siunitx, acronym
% =============================================================================

\chapter{Grundlagen}
\label{chap:grundlagen}

Dieses Kapitel legt das theoretische Fundament für die in den folgenden Kapiteln
beschriebene Entwicklung. Es gliedert sich in fünf Abschnitte: Zunächst werden die
physikalischen Grundlagen der LoRa"=Modulationstechnologie erläutert, gefolgt von
einer Darstellung des LoRaWAN"=Protokollstacks und seiner Netzwerkarchitektur.
Anschließend werden Sicherheitsmechanismen, Energieeffizienz sowie Methoden der
RF"=Planung und Signalausbreitung behandelt. Abschließend wird ein kurzer Überblick
über die Normierungsanforderungen im Bereich Schutzsensorik gegeben.


% -----------------------------------------------------------------------------
\section{LoRa -- Physical Layer}
\label{sec:lora-physical}
% -----------------------------------------------------------------------------

LoRa (\textit{Long Range}) ist eine proprietäre Modulationstechnologie des Halbleiterherstellers
Semtech und bildet die physikalische Grundlage des LoRaWAN"=Protokollstacks. Es ist wichtig,
zwischen LoRa und LoRaWAN zu unterscheiden: LoRa bezeichnet ausschließlich die
Modulationstechnik, während LoRaWAN den darüber liegenden Netzwerkstandard beschreibt,
der von der LoRa Alliance definiert und gepflegt wird~\cite{semtech_lora_basics}.

\subsection{Chirp Spread Spectrum}
\label{subsec:css}

LoRa basiert auf dem Verfahren der \textit{Chirp Spread Spectrum}"=Modulation (CSS).
Dabei wird das zu übertragende Signal auf ein breiteres Frequenzband gespreizt, indem
sogenannte Chirps -- Signale mit kontinuierlich steigender oder fallender Frequenz --
als Trägersymbole verwendet. Dieses Verfahren verleiht LoRa eine hohe Robustheit
gegenüber Mehrwegeausbreitung, Schmalbandstörungen und dem Doppler"=Effekt, wie er
bei mobilen Endgeräten auftritt~\cite{augustin2016study}.

Gegenüber klassischen Schmalband"=Modulationsverfahren wie \ac{FSK} besitzt CSS einen
wesentlichen Vorteil: Das Signal kann noch korrekt demoduliert werden, selbst wenn der
Empfangspegel deutlich unterhalb des Rauschpegels liegt. Das resultierende
Signal"=zu"=Rausch"=Verhältnis (\ac{SNR}) kann dabei negative Werte von bis zu
\SI{-20}{\dB} annehmen. Mit \ac{FSK} lassen sich Datenraten von bis zu \SI{50}{\kbps}
erzielen; LoRa liefert demgegenüber typische Nutzdatenraten zwischen \SI{300}{\bps}
und \SI{11}{\kbps}~\cite{semtech_lora_basics}.

\subsection{Spreading Factor, Bandbreite und Coding Rate}
\label{subsec:sf-bw-cr}

Die wichtigsten Parameter zur Steuerung des LoRa"=Signals sind der Spreading Factor
(\ac{SF}), die Bandbreite (\ac{BW}) sowie die Coding Rate (\ac{CR}).

Der \ac{SF} bestimmt, wie stark das Signal gespreizt wird, und kann ganzzahlige Werte
zwischen 7 und 12 annehmen. Ein höherer \ac{SF} erhöht die Empfindlichkeit des
Empfängers und damit die erreichbare Reichweite, verlängert jedoch gleichzeitig die
\textit{Time on Air} (Luftzeit) eines Paketes erheblich und erhöht den
Energieverbrauch pro Übertragung. Tabelle~\ref{tab:sf_vergleich} fasst die
Auswirkungen der verschiedenen Spreading Factors zusammen.

\begin{table}[htbp]
  \centering
  \caption{Vergleich der LoRa Spreading Factors bei \SI{125}{\kHz} Bandbreite}
  \label{tab:sf_vergleich}
  \begin{tabular}{@{}lrrrl@{}}
    \toprule
    \textbf{SF} & \textbf{Datenrate} & \textbf{Empfindlichkeit} &
    \textbf{Time on Air} & \textbf{Typischer Einsatz} \\
                & [\si{\bps}]        & [\si{\dBm}]              &
    (20\,Byte) [ms] & \\
    \midrule
    SF7  & 5470 & $-123$ &  41 & Kurze Distanz, Energie sparen \\
    SF8  & 3125 & $-126$ &  72 & Ausgewogener Betrieb \\
    SF9  & 1760 & $-129$ & 144 & Mittlere Distanz \\
    SF10 &  980 & $-132$ & 289 & Grenzbereich Stadtgebiet \\
    SF11 &  440 & $-134$ & 578 & Ländliche Umgebung \\
    SF12 &  250 & $-137$ & 991 & Maximale Reichweite \\
    \bottomrule
  \end{tabular}
\end{table}

Als Faustregel gilt: SF7 zusammen mit einer Bandbreite von \SI{125}{\kHz}
(\texttt{SF7BW125}) stellt den optimalen Ausgangspunkt für die meisten Anwendungen
dar, da die Luftzeit und damit der Energieverbrauch minimiert werden~\cite{semtech_lora_basics}.
Die Coding Rate beschreibt den Anteil der redundanten Bits zur Fehlerkorrektur und
nimmt Werte von 4/5 bis 4/8 an.

\subsection{Frequenzbänder und regulatorische Rahmenbedingungen}
\label{subsec:frequenzbaender}

LoRa nutzt lizenzfreie ISM"=Bänder (\textit{Industrial, Scientific and Medical}).
In Europa werden die Frequenzbänder \SI{433}{\MHz} und \SI{868}{\MHz} verwendet,
in Nordamerika \SI{915}{\MHz}. Da die Frequenzzuweisungen regional verschieden sind,
können Geräte, die für den europäischen Markt entwickelt wurden, nicht ohne
Anpassungen in anderen Regionen betrieben werden~\cite{lorawan_spec}.

Für das \SI{868}{\MHz}"=Band in Europa gilt eine Duty"=Cycle"=Begrenzung: Endgeräte
dürfen je nach Subband 0{,}1\,\%, 1\,\% oder 10\,\% der Zeit senden. Die
10\,\%"=Ausnahme gilt ausschließlich für Gateways~\cite{lorawan_spec}. In der Praxis
bedeutet eine 1\,\%"=Beschränkung beispielsweise, dass ein Gerät nach einer
Sendedauer von \SI{1}{\s} mindestens \SI{99}{\s} warten muss, bevor es erneut auf
demselben Kanal sendet. Für zeitkritische Warn"= und Alarmsysteme ist diese
Einschränkung bei der Systemauslegung zwingend zu berücksichtigen.

Die maximale Paketgröße in Europa beträgt 222\,Bytes bei hoher Datenrate und
51\,Bytes bei niedriger Datenrate. Dabei addiert LoRaWAN stets einen
13"=Byte"=Overhead zum Nutzdaten"=Payload~\cite{lorawan_spec}.


% -----------------------------------------------------------------------------
\section{LoRaWAN -- Network Layer}
\label{sec:lorawan-network}
% -----------------------------------------------------------------------------

LoRaWAN definiert den Netzwerkstandard, der auf der LoRa"=Physikschicht aufsetzt.
Die Spezifikation TS001 wird durch das Begleitdokument RP002 ergänzt, das
regionsspezifische Parameter festlegt~\cite{lorawan_spec}.

\subsection{Netzwerkarchitektur}
\label{subsec:netzwerkarchitektur}

Eine LoRaWAN"=Infrastruktur besteht aus drei logischen Ebenen: Endgeräten
(\textit{End Devices}), Gateways und Servern. Die Server"=Seite gliedert sich
weiter in \textit{Network Server}, \textit{Join Server} und
\textit{Application Server}.

\textbf{Endgeräte} bestehen aus einem Mikrocontroller, auf dem der LoRaWAN"=Stack
läuft, einer LoRa"=Funkeinheit, einer Antenne, einer Stromversorgung sowie
anwendungsspezifischen Peripheriegeräten wie Gassensoren. Die Softwareschichten
umfassen Treiber, Middleware und die eigentliche Applikation.

\textbf{Gateways} empfangen Pakete der Endgeräte und leiten diese über eine
IP"=Backhaul"=Verbindung (Ethernet oder Mobilfunk) an den Network Server weiter.
Gateways sind zustandslos (\textit{stateless}) -- sie speichern keine
sitzungsspezifischen Informationen und unterscheiden nicht zwischen verschiedenen
Endgeräten~\cite{lorawan_spec}. In Europa werden je nach Band insgesamt 10 Kanäle
definiert, von denen drei obligatorisch von allen Geräten unterstützt werden müssen.
Bei Systemen für den Außeneinsatz sind Outdoor"=Gateways erforderlich, die aufgrund
wasserdichter Gehäuse und Hohlraumfilter teurer sind als Indoor"=Varianten. Als
Faustegel empfiehlt sich für Indoor"=Installationen die Platzierung des Gateways in
der Nähe von Fensterfronten, um den Signaldurchgang zu verbessern.

Der \textbf{Network Server} ist der zentrale Knotenpunkt des LoRaWAN"=Netzwerks. Er
übernimmt Nachrichtenkonsolidierung (falls ein Paket von mehreren Gateways empfangen
wurde), Routing, Netzwerksteuerung sowie die Überwachung von Gateways und
Endgeräten~\cite{lorawan_spec}. Das \textit{Network Reference Model} (NRM) definiert
die standardisierten Schnittstellen zwischen diesen Elementen.

Der \textbf{Join Server} verwaltet die kryptografischen Wurzelschlüssel und wird
insbesondere beim \ac{OTAA}"=Aktivierungsverfahren benötigt. Er muss in einer
vertrauenswürdigen, gesicherten Umgebung betrieben werden. Um Herstellerabhängigkeiten
(\textit{Vendor Lock"=in}) zu vermeiden, sollte darauf geachtet werden, dass der Join
Server netzwerkunabhängig betrieben werden kann~\cite{semtech_lora_basics}.

LoRaWAN"=Netzwerke verwenden das Aloha"=Zugriffsverfahren: Endgeräte senden, wann
immer sie Daten zu übertragen haben, ohne vorherige Kanalprüfung. Gateways hören
permanent auf allen Kanälen. Kollisionen werden durch das Spread"=Spectrum"=Verfahren
und den Einsatz mehrerer Kanäle minimiert, aber nicht vollständig ausgeschlossen.

\subsection{Geräteklassen}
\label{subsec:geraeteklassen}

Die LoRaWAN"=Spezifikation definiert drei Geräteklassen (A, B und C), die festlegen,
wann ein Endgerät Downlink"=Nachrichten empfangen kann. Die Klasse kann während des
Betriebs gewechselt werden.

\textbf{Klasse A} ist die energieeffizienteste und am weitesten verbreitete Klasse.
Nach jeder Uplink"=Übertragung öffnet das Gerät zwei kurze Empfangsfenster
(RX1 und RX2). Außerhalb dieser Fenster befindet sich das Gerät im Schlafmodus.
Downlink"=Nachrichten können vom Network Server ausschließlich innerhalb dieser
Fenster gesendet werden.

\textbf{Klasse B} ermöglicht periodischen Empfang durch zeitgesteuerte Empfangsfenster,
die mittels Gateway"=Beacons synchronisiert werden. Dies erhöht die Reaktionsfähigkeit
auf Kosten eines höheren Energieverbrauchs.

\textbf{Klasse C} ermöglicht kontinuierlichen Empfang, indem das Endgerät nur während
des Sendens nicht empfängt. Diese Klasse ist für batteriebetriebene Geräte ungeeignet
und wird typischerweise bei netzgespeisten Aktoren eingesetzt.

\subsection{Aktivierungsverfahren: OTAA und ABP}
\label{subsec:aktivierung}

Für die Aufnahme eines Endgeräts in ein LoRaWAN"=Netzwerk existieren zwei Verfahren:
\ac{OTAA} und \ac{ABP}.

Bei \ac{ABP} werden alle benötigten Parameter -- \texttt{DevEUI}, \texttt{DevAddr}
sowie die Sitzungsschlüssel -- vorab im Gerät fest einprogrammiert. Das Gerät kann
sofort ohne Join"=Prozedur senden. Dieses Verfahren ist weniger flexibel und gilt aus
Sicherheitsperspektive als suboptimal, da keine dynamische Schlüsselerneuerung
stattfindet.

Bei \ac{OTAA} initiiert das Gerät beim ersten Start eine \textit{Join Request},
woraufhin der Network"= bzw. Join"=Server einen \textit{Join Accept} zurücksendet
und dabei die Sitzungsschlüssel dynamisch ableitet. \ac{OTAA} ist das empfohlene
Verfahren~\cite{lorawan_spec}. Als \textit{Best Practice} empfiehlt sich, Geräte,
die per ABP provisioniert werden, parallel mit allen für OTAA notwendigen Parametern
(\texttt{DevEUI}, \texttt{JoinEUI} und Root Keys) auszustatten. Alle Geräte müssen
einen persistenten Zählerstand (Frame Counter) auch über Stromausfälle und Resets
hinaus erhalten.

\subsection{Adaptive Data Rate}
\label{subsec:adr}

\ac{ADR} ist ein Mechanismus des Network Servers zur automatischen Anpassung von
Datenrate und Sendeleistung eines Endgeräts. Endgeräte in der Nähe eines Gateways
werden auf höhere Datenraten (niedrigere SF"=Werte) umgeschaltet, was die Luftzeit
verkürzt und die Netzwerkkapazität erhöht. Geräte in größerer Entfernung werden mit
niedrigeren Datenraten betrieben, um ausreichende Verbindungsqualität sicherzustellen.
Wird ein weiteres Gateway zum Netzwerk hinzugefügt, aktualisieren die Endgeräte ihren
SF automatisch~\cite{semtech_lora_basics}. \ac{ADR} ist ausschließlich für stationäre
Geräte geeignet -- für mobile Anwendungen wie das hier betrachtete Schutzsensorsystem
muss ADR deaktiviert oder durch ein geeignetes mobilitätsbewusstes Verfahren ersetzt
werden.

\subsection{Öffentliche vs. private Netzwerke}
\label{subsec:netzwerktypen}

Bei der Wahl der Netzwerkinfrastruktur sind drei Faktoren maßgeblich: Kosten,
Zugangskontrolle und Datenschutz~\cite{semtech_lora_basics}.

Öffentliche Netze (z.B. The Things Network) sind für kleine Gerätezahlen und
Prototyping geeignet, erfordern jedoch eine bestehende Abdeckung und bieten keine
Möglichkeit zur Optimierung durch eigene Gateways. Private Netzwerke eignen sich
für große Gerätezahlen und sicherheitskritische Anwendungen, erfordern jedoch die
Installation und den Betrieb eigener Gateways und Server"=Infrastruktur. Für das
in dieser Arbeit betrachtete Off"=Grid"=Szenario -- den Einsatz in abgelegenen
Gefahrgebieten ohne bestehende Infrastruktur -- ist ein privates, autarkes Netzwerk
die einzig praktikable Option. Als Network"=Server"=Software kommt dabei \textbf{ChirpStack}
zum Einsatz, da dieser vollständig lokal auf einem Raspberry Pi betrieben werden kann
und keine Internetverbindung erfordert. Semtech bietet für Prototyping"=Zwecke zudem
einen kostenlosen Network Server für bis zu drei Gateways und zehn Endgeräte an.


% -----------------------------------------------------------------------------
\section{Sicherheitsmechanismen in LoRaWAN}
\label{sec:sicherheit}
% -----------------------------------------------------------------------------

LoRaWAN setzt auf AES"=128"=Verschlüsselung auf zwei Ebenen: Der \textit{Network
Session Key} (NwkSKey) sichert die Kommunikation zwischen Endgerät und Network
Server, während der \textit{Application Session Key} (AppSKey) die Ende"=zu"=Ende"=Verschlüsselung
der Nutzdaten zwischen Endgerät und Application Server gewährleistet~\cite{lorawan_spec}.

Für den sicheren Betrieb in industriellen und sicherheitskritischen Umgebungen empfiehlt
sich der Einsatz von \textit{Secure Elements} (kryptografische Chips). Diese
Spezialchips speichern die kryptografischen Schlüssel in schreibgeschützten Registern,
beschleunigen Ver"= und Entschlüsselung ohne CPU"=Belastung und ermöglichen die
Verifikation der Firmware"=Authentizität -- wodurch das Risiko eines unbefugten
Firmware"=Overwrite durch einen Angreifer minimiert wird~\cite{semtech_lora_basics}.

Zur Abwehr von Replay"=Angriffen sollten 32"=Bit"=Rahmenzähler für Uplink und Downlink
aktiviert werden. Darüber hinaus kann eine feste Payload"=Länge dabei helfen, ereignisgesteuerte
Aktivitäten vor Angreifern zu verschleiern, die aus variierenden Paketlängen auf den
Systemzustand schließen könnten. Die \textit{Out"=of"=Band Key Provisioning}"=Methode
-- etwa über BLE oder NFC -- stellt sicher, dass Produktionsmitarbeiter keinen Zugriff
auf die Geräteschlüssel erhalten.


% -----------------------------------------------------------------------------
\section{Energieeffizienz in eingebetteten LoRaWAN"=Systemen}
\label{sec:energieeffizienz}
% -----------------------------------------------------------------------------

Energieeffizienz ist eine der zentralen Anforderungen an das in dieser Arbeit
entwickelte System. Der Großteil des Energieverbrauchs entfällt auf den Sendevorgang:
Typischerweise verbraucht das Senden etwa zehnmal mehr Energie als der
Empfangsbetrieb~\cite{semtech_lora_basics}. Für den STM32WLE5 beträgt die Stromaufnahme
im Sendebetrieb bei maximaler Ausgangsleistung von \SI{22}{\dBm} (bei $V_\text{DDPA} =
\SI{3.3}{\V}$) bis zu \SI{87}{\mA}, im Empfangsbetrieb typischerweise \SI{5.5}{\mA},
während der Deep"=Stop"=Modus lediglich \SI{1.08}{\uA} benötigt~\cite{stm32wl_ds}.

\subsection{Sleep"=Modi und Power Management}
\label{subsec:sleep}

Zephyr RTOS bietet ein umfassendes Power"=Management"=Framework, das es ermöglicht,
den Mikrocontroller zwischen Messzyklen in einen energiesparenden Schlafzustand zu
versetzen. Das \textit{Tickless Idle}"=Konzept deaktiviert den System"=Tick"=Timer
in Ruhephasen vollständig; einzelne Peripheriekomponenten können über
\texttt{PM\_DEVICE\_STATE\_SUSPEND} individuell abgeschaltet werden. Der STM32WLE5
unterstützt mehrere Schlafmodi mit unterschiedlichen Aufwachzeiten und
Stromverbräuchen, vom einfachen Sleep über Stop"=Modi bis hin zum Standby.

\subsection{Einfluss des Spreading Factors auf den Energieverbrauch}
\label{subsec:sf-energie}

Die Wahl des Spreading Factors hat direkten Einfluss auf den Energieverbrauch pro
übertragenem Byte. SF12 erzielt zwar die größte Reichweite, verlängert jedoch die
Luftzeit und damit die Zeit, in der der Sender aktiv ist, um etwa das 24"=fache
gegenüber SF7 (vgl. Tabelle~\ref{tab:sf_vergleich}). Die energieoptimale Strategie
besteht darin, den niedrigstmöglichen SF zu verwenden, der noch eine zuverlässige
Übertragung gewährleistet.

\subsection{Strategien zur Datenmengereduktion}
\label{subsec:datenreduktion}

Neben der Wahl des Spreading Factors lässt sich der Energieverbrauch durch die
Reduktion der Sendefrequenz und der Datenmenge weiter minimieren. Anstatt Rohdaten
kontinuierlich zu übertragen, können aggregierte Werte -- beispielsweise Minimum,
Durchschnitt und Maximum über ein Zeitfenster von fünf Minuten -- gesendet werden.
Alternativ können Übertragungen ereignisgesteuert erfolgen, etwa wenn ein Messwert
einen definierten Schwellwert überschreitet oder eine Bewegung detektiert wird.
\textit{Firmware Updates Over the Air} (FUOTA) ermöglicht darüber hinaus die
Aktualisierung der Gerätesoftware ohne physischen Zugang~\cite{semtech_lora_basics}.


% -----------------------------------------------------------------------------
\section{RF"=Planung und Signalausbreitung}
\label{sec:rf-planung}
% -----------------------------------------------------------------------------

RF"=Planung (\textit{Radio Frequency Planning}) bezeichnet die systematische
Vorgehensweise zur Vorhersage der Funkabdeckung eines drahtlosen Systems.
Sie bildet die methodische Grundlage für den in Kapitel~\ref{chap:evaluation}
durchgeführten Vergleich zwischen theoretisch vorhergesagter und praktisch
gemessener Reichweite.

\subsection{Link Budget}
\label{subsec:link-budget}

Das Link Budget beschreibt die Gesamtbilanz aller Signalgewinne und -verluste
entlang der Übertragungsstrecke. Das \textit{Maximum Allowable Path Loss} (MAPL)
ergibt sich aus:

\begin{equation}
  \text{MAPL} = P_\text{TX} + G_\text{TX} + G_\text{RX} - L_\text{cable}
                - \text{NF} - \text{SNR}_\text{min} - \text{Margin}
  \label{eq:mapl}
\end{equation}

wobei $P_\text{TX}$ die Sendeleistung in \si{dBm}, $G_\text{TX}$ und $G_\text{RX}$
die Antennengewinne in \si{dBi}, NF das Rauschen des Empfängers und
$\text{SNR}_\text{min}$ das minimale empfängerseitig tolerierte Signal"=zu"=Rausch"=Verhältnis
beschreiben. Als maximaler Antennengewinn für eine einfache Dipolantenne gilt
$G = \SI{2.15}{dBi}$. Ein Sicherheitsabstand von 5 bis \SI{10}{\dB} wird empfohlen,
um Schwundeffekte und Hindernisse zu kompensieren~\cite{semtech_lora_basics}.
Uplink und Downlink"=Budget sind aufgrund unterschiedlicher Antennenhöhen in der Regel
nicht symmetrisch.

\subsection{Freiraum"=Pfadverlust und Ausbreitungsmodelle}
\label{subsec:pfadverlust}

Der Freiraum"=Pfadverlust (\textit{Free Space Path Loss}, FSPL) beschreibt den
Signalverlust in einer hindernisfreien Umgebung und berechnet sich nach der
Friis"=Gleichung:

\begin{equation}
  \text{FSPL} \, [\si{\dB}] = 20 \log_{10}(d) + 20 \log_{10}(f)
                              + 20 \log_{10}\!\left(\frac{4\pi}{c}\right)
  \label{eq:fspl}
\end{equation}

mit der Distanz $d$ in Metern und der Frequenz $f$ in Hertz. Bei Verdopplung der
Entfernung steigt der Pfadverlust im Freiraummodell um \SI{6}{\dB}; das
\textit{Two"=Ray Ground Reflection Model} ergibt bei Bodennähe sogar einen Verlust
von \SI{12}{\dB} bei doppelter Distanz~\cite{dobkin2012rf}. Für sehr große Distanzen
ergibt sich für das \SI{868}{\MHz}"=Band bei beispielsweise \SI{800}{\km} ein
theoretischer FSPL von \SI{-150}{\dB}.

Für realitätsnähere Vorhersagen werden empirische Ausbreitungsmodelle verwendet.
Tabelle~\ref{tab:ausbreitungsmodelle} gibt einen Überblick der gängigsten Modelle.

\begin{table}[htbp]
  \centering
  \caption{Übersicht der Ausbreitungsmodelle und deren Einsatzbereich}
  \label{tab:ausbreitungsmodelle}
  \begin{tabular}{@{}lll@{}}
    \toprule
    \textbf{Modell} & \textbf{Einsatzbereich} & \textbf{Bemerkung} \\
    \midrule
    Log"=Distance Path Loss & Allgemein & Einfachstes empirisches Modell \\
    Okumura"=Hata            & 150--1500\,MHz, Außenbereich & Meistgenutztes Modell \\
    COST 231                 & 1500--2000\,MHz              & Detaillierte Stadtarchitektur \\
    ITU"=R P.1238            & Innenbereich                 & Für Gebäudeplanung \\
    \bottomrule
  \end{tabular}
\end{table}

Das \textbf{Okumura"=Hata}"=Modell ist für den Frequenzbereich von 150 bis
\SI{1500}{\MHz} das am häufigsten eingesetzte Modell und eignet sich damit direkt für
das \SI{868}{\MHz}"=Band~\cite{haslett2008essentials}. Für Innenräume, wie sie
bei der Messung in Industriegebäuden auftreten, ist das \textbf{ITU"=R P.1238}"=Modell
geeigneter. Bei der Interpretation von Messergebnissen sind zudem
Mehrwegeausbreitung (\textit{Multipath Fading}) und die Fresnel"=Zone zu
berücksichtigen: Eine freie Fresnel"=Zone erfordert, dass das Gelände innerhalb einer
ellipsoidförmigen Zone um die direkte Sichtlinie frei von Hindernissen ist.

\subsection{RF"=Planungsmethodik}
\label{subsec:rf-methodik}

Die empfohlene Methodik zur RF"=Planung folgt einem dreistufigen Prozess~\cite{haslett2008essentials}:

\begin{enumerate}
  \item \textbf{Link"=Budget"=Analyse:} Berechnung des MAPL anhand der Systemparameter
        (Sendeleistung, Antennengewinn, Empfängerempfindlichkeit, Sicherheitsmarge).
  \item \textbf{RF"=Coverage"=Simulation:} Verwendung von Ausbreitungsmodellen und
        digitalem Geländemodell zur kartografischen Darstellung der vorhergesagten Abdeckung.
  \item \textbf{Feldmessungen:} Verifikation der Simulation durch reale Messungen in den
        Zielumgebungen.
\end{enumerate}

Für die Coverage"=Simulation steht mit \textbf{SPLAT!} ein quelloffenes Linux"=Werkzeug
zur Verfügung, das das Longley"=Rice"=Modell mit \ac{SRTM}"=Geländedaten kombiniert.
Die gemessenen RSSI"=Werte können anschließend mit Python und der Bibliothek Folium
als georeferenzierte Heatmap über OpenStreetMap visualisiert und mit der
SPLAT!"=Simulation überlagert werden.

\subsection{Mikrozellen und Makrozellen}
\label{subsec:zellen}

In der Netzplanung wird zwischen \textit{Mikrozellen} -- bei denen die Antenne
unterhalb der Dachlinie installiert wird -- und \textit{Makrozellen} mit Antennen
oberhalb der Dachlinie unterschieden. Diese Unterscheidung beeinflusst das
gewählte Ausbreitungsmodell und die erreichbaren Abdeckungsradien~\cite{walter2006radio}.
Für Indoor"=Installationen gilt als Faustregel: Wenn die Innenwände nicht aus Beton oder
Mauerwerk bestehen, kann ein einzelnes Gateway auf der gleichen Etage wie die
Endgeräte ausreichende Abdeckung liefern.


% -----------------------------------------------------------------------------
\section{Normierung und Anforderungen im Bereich Schutzsensorik}
\label{sec:normierung}
% -----------------------------------------------------------------------------

Der Einsatz mobiler Gasmesssysteme in industriellen Umgebungen unterliegt einer
Reihe normativer Vorgaben. Relevant sind insbesondere die Regelungen der
\ac{DGUV} sowie die ATEX"=Richtlinie 2014/34/EU für Geräte in
explosionsgefährdeten Bereichen. Die EN"=60079"=Normenreihe definiert die Anforderungen
an elektrische Betriebsmittel in explosionsgefährdeten Bereichen hinsichtlich
Zündschutzarten, Gehäuseschutz (IP"=Klassen) und Temperaturklassen. Diese
Anforderungen sind bei der Hardwareauswahl und dem mechanischen Design des Prototyps
zu berücksichtigen und werden in Kapitel~\ref{chap:hardware} eingehend behandelt.


% -----------------------------------------------------------------------------
% Abkürzungsverzeichnis-Einträge (für acronym-Paket)
% In Präambel definieren:
% \acrodef{LoRa}{Long Range}
% \acrodef{SF}{Spreading Factor}
% \acrodef{BW}{Bandwidth}
% \acrodef{CR}{Coding Rate}
% \acrodef{SNR}{Signal-to-Noise Ratio}
% \acrodef{FSK}{Frequency Shift Keying}
% \acrodef{ADR}{Adaptive Data Rate}
% \acrodef{OTAA}{Over-the-Air Activation}
% \acrodef{ABP}{Activation by Personalization}
% \acrodef{MAPL}{Maximum Allowable Path Loss}
% \acrodef{FSPL}{Free Space Path Loss}
% \acrodef{SRTM}{Shuttle Radar Topography Mission}
% \acrodef{DGUV}{Deutsche Gesetzliche Unfallversicherung}
% -----------------------------------------------------------------------------